% Part: first-order-logic
% Chapter: syntax
% Section: intro-syntax

\documentclass[../../../include/open-logic-section]{subfiles}

\begin{document}

\olfileid{fol}{syn}{itx}

\olsection{परिचय}

प्रथम-क्रम तर्कशास्त्राची उपपत्ती आणि अधिउपपत्ती विकसित करण्यासाठी,
आपण प्रथम त्याच्या व्यंजकांची विन्यासमीमांसा आणि चिन्हार्थमीमांसा
परिभाषित केली पाहिजे. प्रथम-क्रम तर्कशास्त्राची व्यंजके म्हणजे पदे आणि
!!{formula}s. पदे !!{variable}s, !!{constant}s आणि !!{function}s
यांपासून तयार होतात. !!^{formula}s मात्र !!{predicate}s आणि पदे
यांपासून तयार होतात (यांतून सर्वांत लहान, ``आण्विक'' !!{formula}s
तयार होतात); मग तार्किक संयोजके आणि संख्यापक वापरून आण्विक
!!{formula}s पासून अधिक जटिल सूत्रे तयार करता येतात. रचनानियम
मांडण्याचे अनेक निरनिराळे मार्ग आहेत; आपण त्यांपैकी केवळ एक शक्य
मार्ग देतो. इतर पद्धती वेगळी चिन्हे निवडतील, संयोजकांचे वेगळे संच
आदिम म्हणून निवडतील आणि कंसांचा वेगळ्या रीतीने वापर करतील (किंवा
तथाकथित पोलिश संकेतपद्धतीप्रमाणे कंस अजिबात वापरणार नाहीत). तरी सर्व
दृष्टिकोनांत सामाईक बाब ही की रचनानियम पदांचा आणि !!{formula}s चा
संच \emph{विगमनाने} परिभाषित करतात. हे योग्य रीतीने केल्यास,
रचनानियमांनुसार प्रत्येक व्यंजक मूलतः एकाच रीतीने तयार होऊ शकते.
त्यामुळे तयार होणारी व्यंजके \emph{एकमेव रीतीने वाचनीय} असतात. परिणामी,
त्यांची ही विगमनाधारित व्याख्या आपल्याला त्याच पद्धतीने---विगमनाधारित
व्याख्येने---त्या व्यंजकांना अर्थ देता येतो, याची खात्री देते.

\end{document}
